\clearpage
\section{GNN Expressiveness}

\subsection{Effect of Depth on Expressiveness}

For the two red nodes in the PDF, determine the minimum number of simplified
sum-aggregation message-passing layers required to produce different
embeddings.

\begin{solution}
Let $v_L$ and $v_R$ denote the red nodes in the left and right graphs.  Since
the update includes both the node itself and its neighbors, their embeddings
are
\[
\begin{array}{c|cccc}
  \text{layer }l & 0 & 1 & 2 & 3 \\ \hline
  h_{v_L}^{(l)} & 1 & 2 & 6 & 18 \\
  h_{v_R}^{(l)} & 1 & 2 & 6 & 17
\end{array}
\]
For example, the neighbor of either red node has embedding 4 after one layer.
After two layers, that neighbor has embedding 12 in the left graph but 11 in
the right graph, which reaches the red node in the third layer.  Thus the
minimum required depth is
\[
  \boxed{3\text{ message-passing layers}}.
\]
\end{solution}

\subsection{Random Walk Matrix}

\begin{enumerate}[label=(\roman*)]
  \item Write the random-walk transition matrix $M$ using the node order in
    the graph.
  \item Find the eigenvector satisfying $\vect{r}=M\vect{r}$, normalize it
    using the $\ell_2$ norm, and round each entry to two decimal places.
\end{enumerate}

\begin{solution}
The edges are
$\{1,2\},\{1,4\},\{2,4\},\{3,4\}$, so the node degrees are
$(2,2,1,3)$.  Using the column-vector convention from the question,
$M_{ij}$ is the probability of moving from node $j$ to node $i$.  Hence
\[
  M=
  \begin{bmatrix}
    0           & \tfrac12 & 0 & \tfrac13 \\
    \tfrac12    & 0        & 0 & \tfrac13 \\
    0           & 0        & 0 & \tfrac13 \\
    \tfrac12    & \tfrac12 & 1 & 0
  \end{bmatrix}.
\]
For an undirected graph, the stationary vector is proportional to the degree
vector.  Therefore an eigenvector satisfying $\vect r=M\vect r$ is
$(2,2,1,3)^\top$.  Its $\ell_2$ norm is $\sqrt{18}$, so the requested
normalized vector is
\[
  \vect r
  =\frac{1}{\sqrt{18}}
    \begin{bmatrix}2\\2\\1\\3\end{bmatrix}
  \approx
    \boxed{[0.47,0.47,0.24,0.71]^\top}.
\]
\end{solution}

\subsection{Relation to Random Walk (i)}

For mean aggregation without a learned transformation or nonlinearity, express
the corresponding random-walk transition matrix using adjacency matrix $A$ and
degree matrix $D$.

\begin{solution}
With the homework's column-distribution convention, the transition matrix is
\[
  \boxed{M=AD^{-1}}.
\]
Indeed, $M_{ju}=A_{ju}/d_u$ is the probability of moving from node $u$ to
node $j$.  Therefore the expected layer-$l$ embedding after one step from
$u$ is
\[
  \sum_j M_{ju}\vect h_j^{(l)}
  =\frac{1}{d_u}\sum_{j\in\mathcal N_u}\vect h_j^{(l)}
  =\vect h_u^{(l+1)}.
\]
Equivalently, if node embeddings are stacked as rows of $H^{(l)}$, the GNN
update is $H^{(l+1)}=D^{-1}AH^{(l)}=M^\top H^{(l)}$.
\end{solution}

\subsection{Relation to Random Walk (ii)}

Repeat 4.3 for the aggregation rule with a one-half skip connection:
\[
  \vect{h}_i^{(l+1)}
  = \frac{1}{2}\vect{h}_i^{(l)}
  + \frac{1}{2|\mathcal{N}_i|}
    \sum_{j\in\mathcal{N}_i}\vect{h}_j^{(l)}.
\]

\begin{solution}
The corresponding lazy-random-walk transition matrix is
\[
  \boxed{M_{\mathrm{skip}}=\frac12 I+\frac12 AD^{-1}}.
\]
The term $\tfrac12 I$ gives a self-transition probability of $1/2$, while the
remaining probability is distributed uniformly among the current node's
neighbors.  In row-stacked embedding form, the update operator is its
transpose, $\tfrac12 I+\tfrac12D^{-1}A$.
\end{solution}

\subsection{Over-Smoothing Effect}

Show that repeated mean aggregation converges as $l\to\infty$ when the graph is
connected and has no bipartite components. Relate the result to the Markov
Convergence Theorem.

\begin{solution}
Let
\[
  P=D^{-1}A,
  \qquad H^{(l+1)}=PH^{(l)}.
\]
Because the graph is connected, the associated random walk is irreducible.
Because the connected graph is non-bipartite, the walk is also aperiodic.
The Markov Convergence Theorem therefore gives
\[
  P^l\longrightarrow \vect{1}\vect{\pi}^\top,
  \qquad
  \pi_i=\frac{d_i}{\sum_j d_j}.
\]
Consequently,
\[
  H^{(l)}=P^lH^{(0)}
  \longrightarrow
  \vect{1}\vect{\pi}^\top H^{(0)}.
\]
Every row of the limiting matrix is the same degree-weighted average of the
initial node embeddings.  Thus every node embedding converges, and all nodes
become indistinguishable under this repeated mean aggregation; this is the
over-smoothing effect.
\end{solution}

\subsection{Learning BFS with a GNN}

Construct a one-dimensional GNN that exactly executes BFS. Specify:
\begin{enumerate}[label=(\alph*)]
  \item a message function,
  \item an aggregation function, and
  \item an update rule.
\end{enumerate}

\begin{solution}
Let $h_v^{(t)}\in\{0,1\}$ indicate whether BFS has reached node $v$ by step
$t$.  The following functions execute BFS exactly:
\begin{align*}
  \text{message:}\quad
    &m_{u\to v}^{(t)}=h_u^{(t)},\\
  \text{aggregation:}\quad
    &a_v^{(t)}=\max_{u\in\mathcal N_v}m_{u\to v}^{(t)},\\
  \text{update:}\quad
    &h_v^{(t+1)}=\max\!\left(h_v^{(t)},a_v^{(t)}\right).
\end{align*}
For an empty neighborhood, define the aggregate to be 0.  The update keeps
every previously reached node at 1 and changes an unreached node to 1 exactly
when it has a reached neighbor.  Hence after $t$ layers, precisely the nodes at
distance at most $t$ from the source have embedding 1.
\end{solution}
